\documentclass[11pt]{article}
\usepackage[margin=1in]{geometry}
\usepackage{amsmath,amssymb,amsthm}
\usepackage{hyperref}
\newtheorem{theorem}{Theorem}
\newtheorem{lemma}[theorem]{Lemma}
\newtheorem{proposition}[theorem]{Proposition}
\newtheorem{corollary}[theorem]{Corollary}
\theoremstyle{definition}
\newtheorem{definition}[theorem]{Definition}
\theoremstyle{remark}
\newtheorem{remark}[theorem]{Remark}
\newcommand{\FD}{\mathrm{FD}}
\newcommand{\Aut}{\operatorname{Aut}}
\newcommand{\Fix}{\operatorname{Fix}}
\newcommand{\N}{\mathbb{N}}
\newcommand{\R}{\mathbb{R}}

\title{Disproof of Conjecture 00000008556:\\
the fixed-point chain of the free distributive lattice has length $n-1$, not $\log p(n)$}
\author{jilint777}
\date{2026-10-04}

\begin{document}
\sloppy
\maketitle

\begin{abstract}
Conjecture 00000008556 asserts three things about the free distributive lattice
$\FD(n)$ on $n$ generators: its automorphism group is the symmetric group $S_n$;
its fixed-point lattice is ``of partition-symmetric type''; and the maximal chain
length of the fixed-point lattice is exactly $\log p(n)$, where $p(n)$ is the
partition number. We refute the third clause, and hence the conjunction. The
automorphisms of $\FD(n)$ are the permutations of the generators, and the
elements they fix are the $n$ threshold functions ``at least $k$ of the $n$
generators''. These form a chain of length $n-1$, so the lengths for
$n=1,2,3$ are $0,1,2$, while $p(1),p(2),p(3)=1,2,3$. The equation
$1=\log_b 2$ forces $b=2$, and then $2=\log_2 3$ is false. So no base $b$ of
the logarithm works, and the natural logarithm fails already at $n=2$. The
argument also covers the bounded lattice (with $0,1$ adjoined), chain length
counted in elements, rounded logarithms, and an ``eventually'' reading. A further
arithmetic fact makes the failure independent of how the chain lengths are
computed: no base $b$ makes both $\log_b 2$ and $\log_b 3$ nonnegative integers.
Everything that the disproof uses is formalized in Lean~4 (core library only):
$\FD(n)$ (as the $\wedge/\vee$-closure of the projections and as the monotone
functions), the $S_n$-action, lattice automorphisms, fixed points, chains, and
the partition numbers.
\end{abstract}

\section{The conjecture and its precise reading}
\begin{quote}
\emph{Definition:} A free distributive lattice is a lattice of ideals.
\emph{Conjecture:} The automorphism group of the free distributive lattice is
the symmetric group; its fixed-point lattice is of partition-symmetric type; and
the maximal chain length of the fixed-point lattice is exactly the logarithm of
the partition number $p(n)$.
\end{quote}
The Chinese version says the same. Here $\FD(n)$ is the free distributive lattice
on $n\ge1$ generators, and $n$ is also the argument of $p(n)$. The conjecture is
the conjunction of
\begin{itemize}
\item[(C1)] $\Aut(\FD(n))=S_n$, acting by permuting the generators;
\item[(C2)] the fixed-point lattice is ``of partition-symmetric type'' (this is
not defined; we do not need it);
\item[(C3)] there is a logarithm $\log=\log_b$ (base $b>0$, $b\ne1$; for example
$b=e$, $2$ or $10$) such that for every $n\ge1$ the maximal chain length
$\ell(n)$ of the fixed-point lattice $\Fix(n)$ equals $\log_b p(n)$.
\end{itemize}
Here $\Fix(n)=\{x\in\FD(n):\varphi(x)=x\text{ for all }\varphi\in\Aut(\FD(n))\}$
is the set of elements fixed by the automorphism group (``\emph{its}
fixed-point lattice''), with the order of $\FD(n)$. The length of a chain
$x_0<x_1<\dots<x_k$ is $k$, the number of steps. We also treat the convention
that counts elements ($k+1$).

We prove $\neg(\mathrm{C1}\wedge\mathrm{C2}\wedge\mathrm{C3})$ for every meaning
of (C2). If (C1) is false the conjunction is false. If (C1) holds, then $\Fix(n)$
is the set of $S_n$-fixed elements, and (C3) fails for it. (C1) is in fact true
(Remark~\ref{rem:aut}), but the proof does not use this.

\paragraph{The definition ``lattice of ideals''.} By Birkhoff's representation
theorem, a finite distributive lattice is the lattice of order ideals (down-sets)
of its poset of join-irreducibles. For $\FD(n)$ this gives the usual descriptions
below: monotone Boolean functions correspond to up-sets of $\{0,1\}^n$, that is,
to the complementary down-sets. Every reading of the definition therefore gives
the same lattice $\FD(n)$, and the definition plays no further role.

\section{Definitions and basic facts}
Points of $\{0,1\}^n$ are identified with subsets $x\subseteq[n]=\{1,\dots,n\}$.
A Boolean function $f:\{0,1\}^n\to\{0,1\}$ is \emph{monotone} if $x\subseteq y$
and $f(x)=1$ imply $f(y)=1$. Functions are ordered pointwise, with $\wedge,\vee$
computed pointwise. The projection $\pi_i$ is $\pi_i(x)=[i\in x]$.

\begin{proposition}\label{prop:fd}
The sublattice of $\{0,1\}^{\{0,1\}^n}$ generated by $\pi_1,\dots,\pi_n$ is the
set of nonconstant monotone functions. It is the free distributive lattice
$\FD(n)$. Adding the constants $0,1$ gives the free bounded distributive
lattice, which consists of all monotone functions. The sizes are
$|\FD(n)|=1,4,18,166$ and $3,6,20,168$ in the bounded case, for $n=1,2,3,4$.
\end{proposition}
\begin{proof}
Projections are monotone with $\pi_i(\emptyset)=0$ and $\pi_i([n])=1$, and these
three properties are preserved by $\wedge$ and $\vee$. Conversely, let $f$ be
monotone and nonconstant. Then $f(\emptyset)=0$, since otherwise monotonicity
gives $f\equiv1$. Hence $f=\bigvee_{x:f(x)=1}\bigwedge_{i\in x}\pi_i$, where
every $x$ in the join is nonempty and the join is nonempty. That this lattice is
free is the classical normal form theorem for distributive lattices (two lattice
polynomials agree in every distributive lattice if and only if they agree on
$\{0,1\}$). The sizes are the Dedekind numbers $3,6,20,168$ minus the two
constants.
\end{proof}

$S_n$ acts on functions by permuting variables: $(\sigma f)(x)=f(\sigma(x))$.
Each $\sigma$ preserves $\wedge$, $\vee$ and $\FD(n)$, so it is a lattice
automorphism.

\begin{proposition}\label{prop:fix}
The $S_n$-fixed elements of $\FD(n)$ are the threshold functions
$T_k(x)=[\,|x|\ge k\,]$ with $1\le k\le n$; in the bounded lattice they are those
with $0\le k\le n+1$. They form the chain $T_n<T_{n-1}<\dots<T_1$, of length
$n-1$ (length $n+1$ in the bounded case).
\end{proposition}
\begin{proof}
If $f$ is fixed by all $\sigma$ then $f(x)$ depends only on $|x|$, because $S_n$
is transitive on the $k$-subsets. Write $f(x)=g(|x|)$. Monotonicity makes $g$
nondecreasing on $\{0,\dots,n\}$, so $g=[\,\cdot\ge k\,]$ for some
$0\le k\le n+1$. The constants are $k=0$ ($f\equiv1$) and $k=n+1$ ($f\equiv0$).
Conversely every $T_k$ is symmetric and monotone. Since $T_{k+1}\le T_k$, with
strict inequality, the fixed set is totally ordered with $n$ (respectively $n+2$)
elements.
\end{proof}

\begin{remark}[Clause (C1) is true]\label{rem:aut}
For a finite distributive lattice $L$, the automorphisms of $L$ correspond to the
automorphisms of its poset $J(L)$ of join-irreducibles (Birkhoff duality). For
$\FD(n)$, $J$ consists of the elements $\bigwedge_{i\in S}\pi_i$ with
$\emptyset\ne S\subsetneq[n]$ (all $S\subseteq[n]$ in the bounded case), ordered
by reverse inclusion. The singletons $\{i\}$ are the minimal elements of
$J\setminus\{\emptyset\}$ under inclusion ($\emptyset$, present only in the
bounded case, is the extreme element and is fixed). So an automorphism $\varphi$
of $J$ permutes the singletons, say $\varphi(\{i\})=\{\sigma(i)\}$. Since
$\{i\}\subseteq S\iff\varphi(\{i\})\subseteq\varphi(S)$ and every $S$ is the union
of the singletons it contains, $\varphi(S)=\sigma(S)$.
Hence
$\Aut(\FD(n))\cong S_n$, acting by permuting the generators. The script
\texttt{verify.py} confirms this by brute-force backtracking over all lattice
automorphisms for $n\le3$, and over all automorphisms of $J$ for $n=4$.
Consequently $\Fix(n)$ is exactly the chain of Proposition~\ref{prop:fix}.
\end{remark}

\paragraph{Partition numbers.} $p(n)$ is the number of nonincreasing sequences of
positive integers with sum $n$: $p(1),\dots,p(5)=1,2,3,5,7$.

\section{The counterexample}
For real $b>0$ with $b\ne1$, real $x>0$ and real $\ell$, by the definition of the
logarithm,
\begin{equation}\label{eq:log}
\ell=\log_b x\iff b^\ell=x .
\end{equation}
Chain lengths are natural numbers, so only powers $b^\ell$ with $\ell\in\N$
occur.

\begin{theorem}\label{thm:main}
Let $\ell(n)$ be the maximal chain length of the $S_n$-fixed lattice of $\FD(n)$.
There is no base $b$ with $\ell(n)=\log_b p(n)$ for all $n\ge1$. The same holds
for the bounded lattice, and for chain length counted in elements. Consequently
conjecture 00000008556 is false.
\end{theorem}
\begin{proof}
By Proposition~\ref{prop:fix}, the lengths at $n=1,2,3$ are $s,s+1,s+2$, where
$s=0$ (free lattice, steps), $s=1$ (free, elements), $s=2$ (bounded, steps) or
$s=3$ (bounded, elements). Also $(p(1),p(2),p(3))=(1,2,3)$. By \eqref{eq:log},
the clause gives $b^s=1$, $b^{s+1}=2$ and $b^{s+2}=3$. The first two give
$b=b^{s+1}/b^s=2$. Then $b^{s+2}=b^{s+1}\cdot b=4\ne3$, a contradiction.
In particular, for the natural logarithm and the free lattice: $\ell(2)=1\ne\ln 2$.
For the conjecture: if (C1) fails, the conjunction fails. If (C1) holds, then
$\Fix(n)$ equals the $S_n$-fixed lattice, and (C3) fails by the above.
\end{proof}

The only algebra used is $b\cdot1=b$, associativity, and $2\cdot2=4\ne3$ for
numbers. In Lean the base $b$ therefore ranges over an arbitrary ``number
system'': any type with an associative unital multiplication and an injective,
multiplicative map from $\N$. Examples are $\N,\mathbb Z,\mathbb Q,\R,\mathbb C$.

\section{Robustness}\label{sec:readings}
\paragraph{Free or bounded lattice; steps or elements.} All four combinations
are covered by Theorem~\ref{thm:main} and by the Lean theorem.

\paragraph{Any base.} The base $b$ is arbitrary: $e$, $2$, $10$, or even
$0<b<1$. The value $b=2$ is forced by $n\le2$, and in fact the clause \emph{does}
hold for $n=1,2$ with $b=2$ (Lean: \texttt{clause\_holds\_n\_le\_2}). The failure is
at $n=3$, where $\log_2 3\approx1.585\ne2$.

\paragraph{Any chain-length convention and any ``fixed-point lattice''.}
\begin{proposition}\label{prop:arith}
For a real $b>0$ and natural numbers $L_2,L_3$, it is impossible that
$L_2=\log_b p(2)$ and $L_3=\log_b p(3)$.
\end{proposition}
\begin{proof}
Otherwise $b^{L_2}=2$ and $b^{L_3}=3$, so $2^{L_3}=b^{L_2L_3}=3^{L_2}$. By parity
$L_3=0$, so $3^{L_2}=1$ and $L_2=0$. Then $b^0=1\ne2$.
\end{proof}
So, whatever is meant by the fixed-point lattice, the clause fails at $n=2,3$,
as long as chain lengths are integers. This includes the fixed lattice of a
single automorphism, which is not what ``its'' (the group's) says. For example
\texttt{verify.py} prints the lengths $2,1$ for $n=2$ and $6,4,2$ for $n=3$ for the
conjugacy classes of $S_2$ and $S_3$ (identity, transposition, $3$-cycle). Proposition~\ref{prop:arith} is proved in Lean as
\texttt{no\_lengths\_fit}.

\paragraph{Rounded logarithms.} ``Exactly'' excludes rounding. Even so,
$\lfloor\cdot\rfloor$, $\lceil\cdot\rceil$ and nearest-integer rounding of
$\log_b p(n)$ for $b\in\{2,e,10\}$ all fail to equal $n-1$ at some $n\le5$
(\texttt{verify.py} prints the first failure). For example
$\lceil\log_2 p(5)\rceil=\lceil\log_2 7\rceil=3\ne4$.

\paragraph{``Eventually'' or asymptotic readings.} For every base $b>1$ the clause
fails for all large $n$, even up to a bounded error. For $0<x<1$ and every $n$,
\[
p(n)\,x^n\le\sum_{m\ge0}p(m)x^m=\prod_{k\ge1}\frac1{1-x^k}=:F(x)<\infty ,
\]
so $p(n)^{1/n}\le x^{-1}F(x)^{1/n}$ and $\limsup_n p(n)^{1/n}\le 1/x$. Letting
$x\to1$ gives $p(n)^{1/n}\to1$, that is, $\log_b p(n)=o(n)$. Hence
$|(n-1)-\log_b p(n)|\to\infty$, and the identity, rounded or not, can hold for
only finitely many $n$. For $0<b<1$ the logarithm is negative for $n\ge2$. This
paragraph is not formalized; the exact statement (Theorem~\ref{thm:main}) is.

\section{Lean formalization}
The project \texttt{lean4/} uses Lean~4.19.0 with the core library only (no
Mathlib). All objects are defined from scratch in \texttt{Main.lean}. A Boolean
function on $n$ variables is a truth table encoded as a natural number
$f<2^{2^n}$: bit $m$ is the value at the point $m$, and coordinate $i$ of $m$ is
bit $i$ of $m$.
\begin{itemize}
\item \texttt{Gen n bnd}: the inductive closure of the projections
\texttt{proj n i} under \texttt{\&\&\&} and \texttt{|||}, with $0$ and $1$ when
\texttt{bnd = true}. \texttt{FDset n bnd}: monotone and, unless bounded,
nonconstant truth tables. \texttt{fd\_spec} proves, for $1\le n\le3$, that these
coincide and equal an explicit list. The proof combines a general soundness
lemma for the closure computation \texttt{closureL} with kernel checks over all
$2^{2^n}$ truth tables. \texttt{fd\_sizes} gives $1,4,18$ and $3,6,20$.
\item \texttt{IsPerm}, \texttt{act}: permutations as duplicate-free lists and
their action on variables. \texttt{Fixed}: the $S_n$-fixed elements.
\texttt{fixed\_iff\_thr}: for $1\le n\le3$, the fixed elements are exactly the
thresholds \texttt{threshold n k}, with $k$ in the admissible range.
\item \texttt{IsAut}: a lattice automorphism (a bijection of $\FD(n)$ preserving
$\wedge,\vee$). \texttt{AutIsSym}: clause (C1). \texttt{FixedAut}: elements fixed
by every automorphism. \texttt{act\_isAut}: each permutation is an automorphism.
\texttt{fixedAut\_iff}: under (C1), \texttt{FixedAut} $=$ \texttt{Fixed}.
\item \texttt{MaxChainLen P k}: some chain in \texttt{P} has $k+1$ elements and
none has more. \texttt{maxChain\_fixed}: lengths $0,1,2$ (free) and $2,3,4$
(bounded) for $n=1,2,3$. The upper bound uses a pigeonhole lemma
(\texttt{nodup\_len\_le}).
\item \texttt{IsPartition}, \texttt{partitionsL}, \texttt{p}: \texttt{mem\_partitionsL}
proves for every $n$ that \texttt{partitionsL n} lists exactly the partitions of
$n$. \texttt{p\_values}: $p(1..5)=1,2,3,5,7$, and the lists for $n\le3$ have no
duplicates.
\item \texttt{NumSys}, \texttt{IsLog R b x y} $:\iff b^y=x$ (equation
\eqref{eq:log}), and \texttt{LogClause R bnd elems}: clause (C3) for the
$S_n$-fixed lattice.
\end{itemize}
Main results:
\begin{itemize}
\item \texttt{logClause\_false}: $\neg$\,\texttt{LogClause R bnd elems}, for all $R$, \texttt{bnd}, \texttt{elems};
\item \texttt{conjecture\_00000008556\_false}: $\neg$\,\texttt{Conjecture R C2 bnd elems}, where
\texttt{Conjecture} is (C1) $\wedge$ (C2) $\wedge$ (C3 for \texttt{FixedAut}), for an
arbitrary predicate \texttt{C2};
\item \texttt{no\_lengths\_fit}: Proposition~\ref{prop:arith};
\item non-vacuity: \texttt{natSys}, \texttt{intSys} are number systems,
\texttt{act\_isAut} shows that \texttt{IsAut} is satisfiable, and
\texttt{clause\_holds\_n\_le\_2} shows that the clause holds for $n\le2$ with $b=2$.
\end{itemize}
There is no \texttt{sorry}, no \texttt{native\_decide} and no added axiom.
\texttt{\#print axioms} reports at most \texttt{propext} and \texttt{Quot.sound}.

\paragraph{Limitations.} Core Lean has no real numbers. The theorems are stated
for an arbitrary \texttt{NumSys}; $\R$ satisfies its axioms, and by
\eqref{eq:log} the real-number clause is an instance of \texttt{LogClause}.
Propositions~\ref{prop:fd} and~\ref{prop:fix} for general $n$, the truth of (C1)
(Remark~\ref{rem:aut}) and the asymptotic paragraph are proved in this report
only; Lean covers $n\le3$, which is all the disproof needs. Clause (C2) is not
formalized; it enters only as an arbitrary predicate.

\section{Independent check}
\texttt{verify.py} (Python 3 standard library) represents a function by its set
of true points and uses a different method. For $n\le4$ and both lattices it
computes $\FD(n)$ as a closure and compares it with the set of monotone functions
for $n\le3$. It computes $\Aut(\FD(n))$ without assumptions: by backtracking
over lattice automorphisms for $n\le3$, and over automorphisms of the
join-irreducible poset for $n=4$. It then finds the fixed points of the whole
group (the thresholds) and the longest chain. It also computes $p(n)$ by Euler's
pentagonal recurrence, cross-checked by direct enumeration, and checks every
reading of Section~\ref{sec:readings}.

\section{Reproduction}
\begin{verbatim}
cd lean4 && lake build
cd .. && python3 verify.py
pdflatex report.tex && pdflatex report.tex
\end{verbatim}
\end{document}
